% ===========================================================================
%  E07 — Acompanhar a aprendizagem e a rede
% ===========================================================================
\chapter{E07 \textperiodcentered{} Acompanhar a aprendizagem e a rede}
\label{ep:e07}

\begin{fichaepico}
\begin{campos}
  \item[Objetivo] Cada nível enxerga o que precisa para agir: o professor, a
    turma e cada estudante; a diretoria, a adoção nas escolas; a rede, o
    agregado --- nunca a criança.
  \item[Para quem] Professor; Coordenação regional; Administrador do contrato.
  \item[Valor] É o que transforma uma missão jogada em aula avaliável, e o que
    a secretaria mostra quando presta contas.
  \item[Estado] \parcial{} --- os resultados da turma estão prontos; os
    relatórios da rede têm ajustes pendentes.
\end{campos}
\end{fichaepico}

\begin{tabularx}{\linewidth}{@{}L{16mm}X L{20mm}@{}}
\toprule
\headrow \thd{ID} & \thd{Feature} & \thd{Estado}\\
\midrule
\id{F07.1} & Resultados da turma & \parcial\\
\zebra \id{F07.2} & Relatório da diretoria & \parcial\\
\id{F07.3} & Análise de produto na rede & \parcial\\
\bottomrule
\end{tabularx}

% ---------------------------------------------------------------------------
\feature{F07.1}{Resultados da turma}{\parcial}

As regras abaixo valem para os resultados de todas as missões. O que é
específico de cada uma está no épico \id{E06}.

\begin{historia}{H07.1.1}{Ler a turma sem ranking}{\pronto}
\hu{Professor}{ver como a turma se saiu na missão}{planejar o debate e a aula
seguinte}
\aceite
\begin{itemize}
  \item Em \enquote{Resultados da turma}, escolho a turma e vejo a folha dela.
  \item Nenhuma folha tem coluna de pontos, de acerto ou ordena crianças.
  \item Sem turma escolhida, a tela explica o que fazer.
  \item A condução ao vivo da aula (largada, Aula 2, debate, revelação) fica na
    mesma tela dos resultados.
\end{itemize}
\end{historia}

\begin{historia}{H07.1.2}{Zerar a turma para repetir a missão}{\pronto}
\hu{Professor}{apagar os dados de uma missão na minha turma}{rodar a missão de
novo com a mesma turma}
\aceite
\begin{itemize}
  \item \enquote{Zerar a turma} pede confirmação, porque não tem volta.
  \item Só aquela missão, naquela turma, é zerada; nenhuma outra missão é
    tocada.
\end{itemize}
\end{historia}

\begin{historia}{H07.1.3}{Ler um estudante sem rotulá-lo}{\pronto}
\hu{Professor}{abrir o relatório de um estudante}{entender o caminho dele na
missão}
\aceite
\begin{itemize}
  \item Existe hoje no unooGPT, na Vila do Lume e no Instanooria (o dossiê).
  \item Uma criança por vez; o relatório nunca afirma omissão a partir de
    ausência de dado, e um valor com uma única medida não recebe número.
  \item A Coordenação e a rede não têm acesso ao relatório de uma criança.
\end{itemize}
\end{historia}

\begin{historia}{H07.1.4}{O texto da criança é lido só pelo professor dela}{\parcial}
\hu{família do estudante}{ter certeza de que o que a criança escreve na
missão é lido só pelo professor da turma}{que ela possa escrever com
sinceridade}
\aceite
\begin{itemize}
  \item O texto livre da criança aparece só para o professor da turma.
  \item A Coordenação, a regional e a rede veem apenas \enquote{entregou / não
    entregou}.
\end{itemize}
\pendencia{no Rastro do Boato e no Manifesto da IA, o texto do \enquote{Meu
percurso} ainda é visível à gestão. A versão nova de Quem mandou isso? segue a
regra: a regra \enquote{Se \ldots, então \ldots} da criança só aparece para quem
dá aula na turma (\id{H06.11.5}). Ver a decisão \id{D03}.}
\end{historia}

% ---------------------------------------------------------------------------
\feature{F07.2}{Relatório da diretoria}{\parcial}

\begin{historia}{H07.2.1}{Acompanhar a adoção na minha regional}{\parcial}
\hu{Coordenação regional}{ver quais escolas estão usando as missões}{apoiar
as escolas que estão atrás}
\aceite
\begin{itemize}
  \item \enquote{Relatório} mostra o consolidado por missão e por escola, com
    \enquote{Turmas abertas}, \enquote{Turmas encerradas} e
    \enquote{Estudantes com registro}.
  \item Só a própria regional, sem seletor.
  \item Sem nome de criança, sem bem-estar, sem ranking.
  \item \enquote{Baixar CSV}.
\end{itemize}
\pendencia{a coluna \enquote{Missão} mostra o código interno (por exemplo
\enquote{exp08ef1b}) em vez do nome da missão.}
\end{historia}

\begin{historia}{H07.2.2}{Comparar as regionais}{\pronto}
\hu{Administrador do contrato}{ver o mesmo relatório para qualquer
regional}{acompanhar a rede inteira}
\aceite
\begin{itemize}
  \item O mesmo relatório, com um seletor de regional.
\end{itemize}
\end{historia}

% ---------------------------------------------------------------------------
\feature{F07.3}{Análise de produto na rede}{\parcial}

\begin{historia}{H07.3.1}{Saber se a missão funciona entre turmas}{\parcial}
\hu{Administrador do contrato}{comparar o envolvimento e a eficácia da missão
entre turmas}{saber onde ela funciona e onde precisa de apoio}
\aceite
\begin{itemize}
  \item \enquote{Análise de produto} mostra, por atividade, Sessões,
    Envolvimento, Tempo de feed, Inversão, \enquote{Achou a busca},
    \enquote{Conteúdo morto} e \enquote{Imagens ok}, com o detalhe de cada
    atividade.
  \item Nenhum nome de estudante.
\end{itemize}
\pendencia{só se chega a esta tela digitando o endereço (não há item de
menu); cobre só o Instanooria; a escola é digitada à mão; e o CSV existe mas
não tem botão na tela.}
\end{historia}

\begin{historia}{H07.3.2}{Ver os números do dia}{\afazer}
\hu{Administrador do contrato}{ver quantas turmas e estudantes usaram a
plataforma hoje}{acompanhar a implantação sem esperar o fim do mês}
\aceiteproposta
\begin{itemize}
  \item Um painel com turmas com missão aberta, estudantes que entraram e
    missões concluídas no dia, por regional.
  \item Nenhum dado individual.
\end{itemize}
\end{historia}
